\section{Drainage Hydrology: Flow Pathways, Regime, and Scaling}

Subsurface drainage reshapes the field and catchment water balance in ways that are now well quantified across observational, tracer, and modelling lines of evidence. Tile networks intercept water that would otherwise recharge shallow groundwater or return as surface runoff and deliver it rapidly to streams, so drained landscapes carry a large and measurable share of streamflow, a distinctive flow regime, and a homogenised, scale-invariant response. This section establishes that physical hydrology---the flow pathways, the modified regime, and its scaling---on which the later water-quality and management findings depend.

\subsection{Tile Drainage as a Dominant Flow Pathway}

Watershed-scale studies consistently establish that tile drains route a large fraction of precipitation and streamflow. In an Ohio watershed, 21\% of precipitation (206~mm) was recovered through tiles, accounting for 41\% of annual and 47\% of monthly discharge \parencite{King2014}, and equivalent-medium modelling of an Iowa watershed attributed 97~mm of tile drainage---13\% of precipitation and 39\% of streamflow \parencite{Thomas2016}. In the Boone River, end-member mixing and \textsc{SWAT} bracketed the tile contribution at roughly 46--54\% of discharge, with \textsc{SWAT} estimating about 53.7\% (Boone) and 65.8\% (Lyons Creek) of annual yield \parencite{Schilling2019}; snowmelt-dominated modelling placed tile flow near 16\% of annual precipitation (about 37\% of yield), with tiling reducing surface runoff and increasing subsurface flow and yield \parencite{Rahman2014}. Field-scale export can be even more tile-dominated: tile discharge accounted for $83\pm20$\% of Ohio water export \parencite{Williams2018}, and most discharge from an Indiana closed depression left through tiles rather than surface flow \parencite{Williams2019}. Across humid settings the drain share commonly falls in the 40--65\% range.

Partitioning between tile and surface pathways varies systematically with hydrologic condition and drainage area. In Iowa, tile drainage averaged 30\% of streamflow during events but 61\% during non-event periods; the event ratio $Q_T/Q$ rose with drainage area while the non-event ratio fell, and annual $Q_T/Q$ increased log-linearly from 0.29 at 6~km$^2$ to 0.39 at 45~km$^2$ ($R^2=0.98$) \parencite{Thomas2016,King2014}.

Tile contributions peak when soils are wet and evapotranspiration low. Ohio tile discharge was strongly seasonal (median monthly summer 0.9~mm versus winter 23.4~mm) \parencite{King2014}, Boone River drainage dominated flow from March through June \parencite{Schilling2019}, and simulations indicated the largest baseflow effect from March to July---over 40\% more baseflow in drained watersheds and about 60\% more in April--June \parencite{Schilling2008}. In cold regions tile flow tracks snowmelt, with a seasonal absence from May to September \parencite{Kompanizare2024}, and water-table position relative to tile depth gates drainage: in a well-drained Iowa block the water table lay below tile depth 72\% of the time, steepening gradients during wet periods \parencite{Logsdon2020}.

\subsection{Modification of the Flow Regime}

\subsubsection{Low Flows and Baseflow}

A robust, repeated finding is that tile drainage elevates low flows and the baseflow fraction. Iowa modelling showed drained landscapes carried far higher baseflow (150~mm, 72\%, at 24~m spacing) than undrained (52~mm, 25\%), with drainage acting primarily on the baseflow portion of the hydrograph \parencite{Schilling2008}. Ohio simulations found tiles greatly increased low flow while decreasing higher flows less, raised annual average daily flow about 5\%, and lowered the probability of the water table entering the root zone from about 50\% to about 5\% \parencite{Rathore2025}. Routing analysis reinforced that tiles increase low flows, decrease intermediate flows, and minimally affect the largest flows \parencite{Sloan2017}. In Latvia, drainage raised mean summer--autumn runoff by 132--186\% and winter low-water runoff by 53--286\%, with annual mean, minimum, and maximum discharge up 46\%, 158\%, and 26\% \parencite{Apsite2017}. Channelized Ohio headwaters showed far fewer zero-discharge days (8 versus 87.5~d~yr$^{-1}$ in natural streams), inverting part of the accepted headwater paradigm \parencite{King2009}.

\subsubsection{Peak Flows and Flood Response}

Effects on the largest peaks are consistently modest, but scale- and condition-dependent for smaller ones. Routing modelling concentrated peak-flow reductions at intermediate scales (100--1000~km$^2$), falling from 26\% at the watershed to 11\% at the basin outlet for one event, with negligible field-scale effect on the largest floods; placement mattered, with draining headwaters producing the greatest reduction \parencite{Sloan2017}. Field-scale Iowa scenarios showed no significant difference in mean annual peak discharge across six soil types, indicating drainage chiefly alters pathway rather than the largest peaks \parencite{Sloan2016}. Ohio simulations found tiles could enhance event peaks for sub-threshold or dry-condition events but reduced them for larger events, cutting streamflow in 23 of the 25 largest events \parencite{Rathore2025}. Under expanded tiling, smaller peaks became more frequent while larger peaks declined, steeper rising limbs implying that reduced large-event peaks arose from reduced volumes rather than altered timing \parencite{Rahman2014}. Distributed modelling similarly found drainage decreases peak flow while increasing recession volume, reproducing observed behaviour with only 21\% of the area tiled \parencite{Fonley2021}.

\subsubsection{Flashiness, Drought, and Long-Term Trends}

Flashiness responds to drainage in a soil- and slope-dependent way. A global meta-analysis of 464 watersheds found overall flashiness (Richards--Baker Index) similar regardless of drainage extent but diverging by setting: in flat watersheds ($<$2\% slope) medium/high drainage raised the index, whereas in sloping watersheds low extent had the highest, with drainage extent correlating positively with agricultural land use and clay content and negatively with slope \parencite{Zhou2025}. Field-scale Iowa simulations showed drainage decreased flashiness for low-conductivity soils and increased it for high-conductivity soils, homogenising response with an optimal spacing near 15~m \parencite{Sloan2016}, while Ohio modelling found overall flashiness decreased with tiles, minimised near 50~m spacing \parencite{Rathore2025}. Across 30 Nordic--Baltic catchments, flashiness related negatively to baseflow index and positively to drained share and intensity, and computing it from hourly rather than daily discharge increased it nearly threefold \parencite{Deelstra2014,Deelstra2014b}. At the regime scale, drainage significantly increased streamflow-drought duration and intensity while reducing frequency, each 10\% increase in tile-drained area lengthening drought duration by 0.03~d, raising intensity 12\%, and lowering frequency 0.10 events~yr$^{-1}$ \parencite{Adelsperger2024}. Indiana trend analysis found increasing low- and medium-flow trends with weakened precipitation linkage in heavily drained areas \parencite{Kumar2009}, although a synthesis of seven Midwestern rivers concluded that inter-annual variability in streamflow and baseflow was controlled primarily by current- and previous-year precipitation \parencite{Baeumler2020}.

\subsection{Scale Dependence, Homogenization, and Travel Times}

Artificial drainage homogenises catchment response and reduces the scale dependence of hydrologic signatures. Non-tiled Iowa watersheds show decreasing area-normalized peak discharge and increasing baseflow ratio with scale, whereas tiled basins behave scale-invariantly; within-basin variability in recession coefficients declines as drained fraction rises, and the tiled/non-tiled distinction disappears above roughly 2200~km$^2$ \parencite{BolandBrien2014}. In Minnesota's Redwood basin, mean travel time fell 41\% (77 to 45~d), year-to-year variability was reduced, and travel time became independent of catchment size above about 200~km$^2$ \parencite{DaneshYazdi2016}, with streamflow dynamics becoming more linear after 1976 despite little change in precipitation \parencite{FoufoulaGeorgiou2015}. Isotope work in a UK catchment found short mean transit times ($<$1~yr, faster in artificial drains) and concluded that drainage homogenised distributed responses \parencite{DimitrovaPetrova2020}, and an Indiana watershed behaved as an effectively homogeneous unit describable by a single exponential travel-time distribution with a 2--3~d mean \parencite{Basu2010}; even so, tile contributions remained detectable beyond 16{,}000~km$^2$ in the Boone River \parencite{Schilling2019}.

Modelling and tracers quantify how drainage shortens groundwater travel times. In Iowa, raising drainage density shortened mean travel time and raised the tile contribution to baseflow to an upper bound of 37\%, while increasing channel incision from 0.3 to 2.7~m lowered mean travel time near-linearly and raised the tile contribution from 30\% to 54\% \parencite{Schilling2012}. In Bear Creek, mean groundwater travel time fell from 5.57 to 1.07~yr at high drainage intensity (versus 181.6~yr presettlement), cutting riparian-buffer residence time to 15.5--36.7~d \parencite{Schilling2015}. Analytical solutions confirm that under tile drainage groundwater levels are lower and decline faster, total groundwater discharge to the river increases, and the tile contribution grows with density, though direct groundwater discharge during recession is reduced \parencite{Liang2015}; recession of heavily tiled watersheds is correspondingly more linear \parencite{Schilling2008}.

Runoff generation is governed by pathway-activation thresholds. In an Indiana closed depression, rainfall alone related linearly to cumulative flow with no threshold, but a clear threshold emerged when antecedent soil-moisture deficit and rainfall were summed; tile flow generally began earlier, peaked earlier, and lasted longer than surface flow \parencite{Williams2019}. Tile celerity there was controlled by antecedent wetness---preferential flow mattering more when dry and stored pre-event water dominating when wetter---with groundwater the primary tile source (mean transit time 245~d) \parencite{Williams2021}. In near-level Canadian vertisols, overland flow generally activated before tile flow, and direct preferential flow into tiles was uncommon \parencite{Kokulan2019}. Across 795 US--Canadian events, 43\% produced both surface and tile flow and 57\% tile flow only, surface flow contributing a median 12\% among dual-pathway events \parencite{Macrae2019}; in Kansas tile-outlet terraces, baseflow contributed 38--51\% of discharge \parencite{Stops2022}.

Tracer and geochemical studies partition tile flow into fast preferential and slower matrix components whose balance shifts with season and wetness. In a French field, dry-condition rainwater transferred directly through macropores whereas saturated conditions were dominated by pre-event soil water; surface and tile drainage transferred 44\% and 20\% of annual effective rainfall, with 65\% of infiltrating water not captured by tiles \parencite{Gaillot2023}, and catchment monitoring found the fast preferential peak contributing about 15\% of water and sediment fluxes \parencite{Grangeon2021}. Ohio master-recession analysis apportioned annual tile discharge into 9\% quick-old, 39\% quick-new, 42\% slow-old, and 10\% slow-new components \parencite{Nazari2021}, preferential flow accounted for a higher fraction in clay than loam \parencite{Pluer2020}, and Swedish work attributed autumn events to preferential flow and spring events to matrix flow \parencite{Salazar2011}. Flow is highly concentrated in time: in Ohio, surface runoff occurred $41\pm42$ and tile drainage $240\pm69$~d~yr$^{-1}$, with 80\% of tile discharge delivered in 21\% of the year \parencite{Williams2018}.

\subsection{Repartitioning of Water Relative to Undrained Land}

Relative to undrained or surface-drained conditions, improved subsurface drainage generally reduces surface runoff, peak flow, sediment, and phosphorus even as it increases nitrate delivery \parencite{FAUSEY1995,SKAGGS1994,THOMAS1995,Gramlich2018}. In Oregon, drainage reduced watershed runoff and sediment yield by about 65\% and 55\%, eliminated baseflow, and lowered peak flow \parencite{ISTOK1983}; Albanian sites recorded surface-runoff reductions of 30--56\% \parencite{Grazhdani1996}; and Minnesota modelling found reduced runoff, soil loss, and sediment yield, especially under no-till \parencite{Maalim2013}. In North Carolina, improved drainage cut surface runoff by 192~mm~yr$^{-1}$ (66\%) while raising total outflow by 10\% and strongly attenuating peak runoff (101 to 28~mm~day$^{-1}$) \parencite{KONYHA1992}.

At the field-water-balance scale, the tile drain, not the soil surface, is usually the principal conduit in humid settings. In Indiana, surface runoff was only 1--10\% of annual precipitation versus 26--45\% for subsurface drainage (simulated balance 7\% runoff, 34\% drain flow, 51\% evapotranspiration), with all ponding generated by saturation excess \parencite{Saadat2020}. Iowa pothole modelling found tile discharge and evapotranspiration accounting for 52.4\% and 47.5\% of precipitation \parencite{Amado2018,Martin2019}, and UK field drains were the dominant pathway transferring runoff and most sediment to the stream, responding in every monitored event \parencite{Deasy2009}. The direction of these shifts is texture-dependent: on sandy loams subsurface drainage minimised runoff and total-phosphorus loss whereas the reverse held on clay loam, phosphorus speciation being mainly a function of soil texture \parencite{Eastman2010}. Mechanistically, good drainage increases infiltration and lowers solute concentrations in surface runoff \parencite{Tong2016,Tao2017}, and by lowering the water table it limits saturation-excess runoff and surface sealing \parencite{Augeard2005}.
